\documentclass[border=4pt]{standalone}
\usepackage[utf8]{inputenc}
\usepackage[T1]{fontenc}
\usepackage[spanish]{babel}
\usepackage{tikz}
\usetikzlibrary{arrows.meta,positioning,shapes.geometric,calc,fit,backgrounds}
\usepackage{amsmath}

\begin{document}
\begin{tikzpicture}[
  font=\small,
  node distance=6mm,
  box/.style={draw, rounded corners=2mm, align=left, text width=12cm, inner sep=6pt},
  arrow/.style={-Latex, thick}
]

\node[box] (s1) {%
\textbf{(1) Universo y marco editorial}\\
Revistas sudamericanas de acceso abierto, indexadas, elegibles según criterios editoriales.
};

\node[box, below=of s1] (s2) {%
\textbf{(2) Selección de revistas (1ra etapa)}\\
Muestreo probabilístico de $J$ revistas desde $J_0$.
};

\node[box, below=of s2] (s3) {%
\textbf{(3) Ventana temporal y listado de candidatos}\\
Último volumen disponible 2024, si no existe 2024 se usa 2023. Enumeración de artículos del volumen auditado.
};

\node[box, below=of s3] (s4) {%
\textbf{(4) Acceso a texto completo y elegibilidad}\\
Descarga PDF, verificación, filtrado por cuantitativo o mixto con inferencia, exclusiones y duplicados.
};

\node[box, below=of s4] (s5) {%
\textbf{(5) Extracción y codificación metodológica}\\
Variables de diseño y reporte, e indicador $y_{ij}$ de falla crítica inferencial (protocolo auditado).
};

\node[box, below=of s5] (s6) {%
\textbf{(6) Control de calidad y base final}\\
Doble codificación (submuestra), Kappa, resolución de discrepancias, consolidación ($n=319$, $J=50$).
};

\draw[arrow] (s1) -- (s2);
\draw[arrow] (s2) -- (s3);
\draw[arrow] (s3) -- (s4);
\draw[arrow] (s4) -- (s5);
\draw[arrow] (s5) -- (s6);

\end{tikzpicture}
\end{document}
